\documentclass[aspectratio=169,11pt]{beamer}
\usepackage[spanish]{babel}
\usepackage[T1]{fontenc}
\usepackage[utf8]{inputenc}
\usepackage{lmodern}
\usepackage{ragged2e}
\usepackage{xcolor}

\usetheme{Madrid}
\setbeamertemplate{navigation symbols}{}
\setbeamertemplate{footline}[frame number]

\definecolor{azul}{RGB}{33,87,146}
\setbeamercolor{structure}{fg=azul}
\setbeamercolor{frametitle}{fg=white,bg=azul}
\setbeamercolor{title}{fg=azul}

\title{Semana 03: Antecedentes Bibliograficos}
\author{Grupo 6}
\date{CIB12 - Aprendizaje de Maquina y Mineria de Datos}

\begin{document}

\begin{frame}
  \titlepage
\end{frame}

\begin{frame}{Tema del Grupo}
\justifying
\textbf{Tema tentativo:}

\vspace{0.5em}

Modelo de aprendizaje profundo con mecanismos de atencion para mejorar la calidad del servicio de guia virtual en servicios de guia virtual con aplicaciones moviles e IoT con senales Wi-Fi/BLE en museos y espacios culturales, 2026.

\vspace{0.8em}

\textbf{Objetivo de la semana 03:}
\begin{itemize}
  \item identificar antecedentes bibliograficos directamente relacionados;
  \item reconocer vacios de la literatura;
  \item sustentar la pertinencia del tema del Grupo 6.
\end{itemize}
\end{frame}

\begin{frame}{Criterios de Seleccion}
\justifying
Se priorizaron trabajos que aportan directamente a:
\begin{itemize}
  \item localizacion indoor con Wi-Fi/BLE;
  \item heterogeneidad de dispositivos;
  \item variaciones temporales del RSSI;
  \item mecanismos de atencion y aprendizaje profundo;
  \item eficiencia para uso en dispositivos moviles;
  \item posibles mejoras para servicios de guia virtual.
\end{itemize}
\end{frame}

\begin{frame}{Antecedentes Centrales}
\justifying
\begin{itemize}
  \item \textbf{Gufran et al. (2023) - STELLAR:} robustez frente a variaciones temporales y heterogeneidad.
  \item \textbf{Tiku et al. (2022) - ANVIL:} mejora la precision ante cambio de smartphone.
  \item \textbf{Singampalli et al. (2024) - xKAN:} interpretabilidad para comprender el impacto del RSSI.
  \item \textbf{Jianhua et al. (2026):} fusion sensorial para precision submétrica.
  \item \textbf{Salamah et al. (2016):} reduccion de costo computacional con PCA.
\end{itemize}
\end{frame}

\begin{frame}{Hallazgos Relevantes}
\justifying
\begin{itemize}
  \item La heterogeneidad de dispositivos afecta fuertemente la precision de localizacion.
  \item Los mecanismos de atencion ayudan a seleccionar senales mas estables y utiles.
  \item La robustez temporal es clave para evitar recalibraciones frecuentes.
  \item La fusion sensorial puede mejorar precision y experiencia de usuario.
  \item La eficiencia computacional sigue siendo necesaria para aplicaciones moviles reales.
\end{itemize}
\end{frame}

\begin{frame}{Sintesis Bibliografica}
\justifying
La literatura revisada muestra una evolucion desde enfoques clasicos de localizacion indoor hacia arquitecturas profundas con atencion, robustez temporal e interes por la heterogeneidad de dispositivos.

\vspace{0.7em}

Los antecedentes seleccionados sostienen que el problema del Grupo 6 es pertinente: todavia existe necesidad de modelos que mejoren la calidad del servicio de guia virtual en museos, manteniendo precision, estabilidad y viabilidad en dispositivos moviles.
\end{frame}

\begin{frame}{Vacios y Pertinencia}
\justifying
\textbf{Vacios detectados:}
\begin{itemize}
  \item dependencia de datasets estaticos;
  \item baja interpretabilidad de varios modelos;
  \item sensibilidad a cambios de dispositivo y del entorno;
  \item poca traduccion directa al contexto de guia virtual en museos.
\end{itemize}

\vspace{0.6em}

\textbf{Pertinencia del tema:}
integrar mecanismos de atencion en un modelo que contribuya a mejorar la calidad del servicio de guia virtual en museos y espacios culturales.
\end{frame}

\begin{frame}{Siguiente Entregable}
\justifying
\textbf{Entrega actual:} antecedentes bibliograficos.

\vspace{0.7em}

\textbf{Siguiente entregable:} dimensiones y diseno.

\vspace{0.7em}

\textbf{Trabajo inmediato del grupo:}
\begin{itemize}
  \item definir variables y dimensiones;
  \item proponer indicadores observables;
  \item aterrizar el diseño preliminar de investigacion.
\end{itemize}
\end{frame}

\end{document}
